\documentclass[12pt,a4paper]{article}
\usepackage[french]{babel}
\usepackage[T1]{fontenc}
\usepackage{amsmath, amssymb}
\usepackage{geometry}
\geometry{top=1cm, bottom=2cm, left=1cm, right=1cm}
\usepackage{multicol}
\usepackage{graphicx}
\usepackage{tcolorbox}
\usepackage{multicol}

\begin{document}

\centering\textbf{Devoir Surveill\'e de Math\'ematiques - Nombres et calculs}\\
\vspace{0.5cm}

\noindent\textbf{Nom :} \underline{\hspace{5cm}} \hfill \textbf{Pr\'enom :} \underline{\hspace{5cm}} \\
\textbf{Classe :} \underline{\hspace{5cm}}\\
\vspace{1cm}

\centering\textit{\textbf{Calculatrice interdite. Toute r\'edaction ou justification pertinente sera valoris\'ee.}}

\section*{Exercice 1 : Appartenance \`a un ensemble (3 points)}

\begin{enumerate}
    \item Compl\'eter par \(\in\), \(\notin\), \(\subset\) ou \(\not\subset\) :
    \begin{tcolorbox}[colframe=black, colback=white, width=\linewidth]
    \begin{itemize}
        \item \(3,5 \ldots \mathbb{N}\)\\[0.2cm]
        \item \(\frac{2}{3} \ldots \mathbb{Q}\)\\[0.2cm]
        \item \(\mathbb{Q} \ldots \mathbb{Z}\)\\[0.2cm]
        \item \(-\sqrt{2} \ldots \mathbb{R}\)\\[0.2cm]
        \item \(-\frac{1}{4} \ldots \mathbb{D}\)\\[0.2cm]
        \item \(\mathbb{N} \ldots \mathbb{R}\)
    \end{itemize}
    \end{tcolorbox}

    \item Donner la nature des nombres suivants :
    \begin{tcolorbox}[colframe=black, colback=white, width=\linewidth]
    \(\frac{25}{5}\in \ldots \), \( -67 \in \ldots \), \( \sqrt{49} \in \ldots \), \\
    \( \frac{17}{23} \in \ldots \), \( 2\pi \in \ldots \), \( \frac{21}{14} \in \ldots \)
    \end{tcolorbox}
\end{enumerate}

\section*{Exercice 2 : D\'ecimaux et d\'emonstration (2,5 points)}
D\'emontrer que les nombres suivants appartiennent \`a \(\mathbb{D}\).
\begin{tcolorbox}[colframe=black, colback=white, width=\linewidth, height=6cm]
\textbf{1.} \(x = 2{,}25\) \\
\textbf{2.} \(y = 2{,}3 \times 10^{-12}\) \\
\textbf{3.} \(z = \frac{19}{50}\)
\end{tcolorbox}

\section*{Exercice 3 : Calculs de fractions (4 points)}
\begin{multicols}{2}
\begin{itemize}
    \item \(A = \frac{4}{21} - \frac{2}{7} + \frac{11}{3}\)
    \item \(B = 5 - 3 \times \frac{4}{11}\)
    \item \(C = (-3) \times \frac{3}{4} + \frac{1}{6}\)
    \item \(D = \frac{-25}{6} \div \frac{10}{3}\)
\end{itemize}
\end{multicols}
\vspace{-0.5cm}
\begin{tcolorbox}[colframe=black, colback=white, width=\linewidth, height=8cm]
\end{tcolorbox}

\section*{Exercice 4 : Intervalles (5 points)}
\begin{enumerate}
    \item \textbf{Donner sous forme d'intervalles ou r\'eunion d'intervalles :}
    \begin{tcolorbox}[colframe=black, colback=white, width=\linewidth, height=4cm]
    \begin{itemize}
        \item les r\'eels strictement positifs : \\
        \item les r\'eels \(\leq -2\) : \\
        \item les r\'eels \(\geq -3\) et \(< 7\) : \\
        \item les r\'eels \(>1\) ou \(\leq -4\) :
    \end{itemize}
    \end{tcolorbox}

    \item \textbf{Repr\'esenter et calculer :} 
    \begin{tcolorbox}[colframe=black, colback=white, width=\linewidth, height=6cm]
    a) \(I = [-3 ; 5],\ J = ]-4 ; 7[\) \\
    b) \(I = ]-\infty ; 2[ ,\ J = ]-1 ; +\infty[\) \\
    c) \(I = ]-3 ; 1[ ,\ J = [1 ; +\infty[\)
    \end{tcolorbox}
\end{enumerate}

\section*{Exercice 5 : Encadrements (1 point)}
\begin{tcolorbox}[colframe=black, colback=white, width=\linewidth, height=4cm]
\textbf{1.} Donner un encadrement d’amplitude \(10^{-1}\) de \(388{,}558\) \\
\textbf{2.} Donner un encadrement d’amplitude \(10^{-2}\) de \(4{,}272\)
\end{tcolorbox}

\section*{Exercice 6 : Variables et algorithme (2 points)}
\noindent On initialise : \(x = 5\), \(y = 4\)
\begin{verbatim}
z \leftarrow 2x - y
y \leftarrow 2y - 3z
x \leftarrow 5z + y - 4x
\end{verbatim}
\begin{tcolorbox}[colframe=black, colback=white, width=\linewidth, height=3cm]
\textbf{Valeurs finales :} \quad \(x = \ldots\) \quad \(y = \ldots\)
\end{tcolorbox}

\section*{Question Bonus : Raisonnement par l’absurde (1 point)}
\begin{tcolorbox}[colframe=black, colback=white, width=\linewidth, height=5cm]
\textbf{Montrer que \(\frac{\pi}{2}\) n’est pas rationnel si \(\pi\) ne l’est pas.}
\end{tcolorbox}

\end{document}
